\section{Geometric robustness: smooth families with transverse folds}
\label{sec:generic-folds}
The optimized threshold does not require the cosine symmetry. This
section proves the same three moment orders for smooth one-parameter
kinetic families. The hypotheses allow unequal root slopes and additional
simple roots, including roots that do not move with the parameter.
The conclusion concerns orders; the sharp cosine coefficients in
\cref{thm:blind-sharp} are not assigned to other families.

\begin{assumption}[A compact fold family]\label{ass:generic-folds}
Let $\Gamma$ be a connected closed wall parametrized by arclength,
let $I$ be a compact parameter interval, and let
$g\in C^4(\Gamma\times I)$. The rates are $k_c(s)=g(s,c)^2$.
The set
\[
 \{(s,c):g(s,c)=g_s(s,c)=0\}
       =\{(s_j,c_j):1\le j\le n_f\}
\]
is finite and nonempty, and every member satisfies
\[
 c_j\in\operatorname{int}I,\qquad
 g_{ss}(s_j,c_j)\ne0,\qquad g_c(s_j,c_j)\ne0.
\]
The sites $s_j$ are pairwise distinct, as are the parameters $c_j$.
The parameter has a probability density $\varpi\in L^\infty(I)$.
In a two-sided neighborhood of each $c_j$, the density is bounded
below almost everywhere by a positive constant.
\end{assumption}

All zeros other than the listed folds are simple in space. In particular,
no realization vanishes identically. Distinct sites and parameters are
explicit sufficient hypotheses: they permit disjoint spatial and
parameter neighborhoods in the proof. We do not assert a simultaneous-fold
extension here. A smooth density positive at every fold is an example
of the sampling condition; positivity away from the folds is unnecessary.

Write $J_c^g(D)=J_{g(\cdot,c)^2}(D)$ using the smooth-test convention
\eqref{eq:J-definition}, and define the predetermined value
\begin{equation}\label{eq:generic-value}
 \mathcal P_q^g(M)=
 \inf_{D\in L^1(\Gamma),\ D\ge0,\ \int_\Gamma D=M}
       \int_I J_c^g(D)^q\varpi(c)\dd c.
\end{equation}
The family and $q>0$ are fixed as the dimensionless budget tends to zero.

\begin{theorem}[Generic optimized moment orders]\label{thm:generic-orders}
Under \cref{ass:generic-folds},
\begin{equation}\label{eq:generic-three}
 \mathcal P_q^g(M)\asymp
 \begin{cases}
 M^{-q/4},&0<q<8/5,\\
 M^{-2/5}[\log(1/M)]^{7/5},&q=8/5,\\
 M^{(2-3q)/7},&q>8/5.
 \end{cases}
\end{equation}
The constants may depend on the family, density, wall, and fixed moment
order, but not on $M$ or on a competing mobility. The lower bounds
hold over all nonnegative integrable fields. Bounded fields with a
strictly positive lower bound attain each upper order.
\end{theorem}

\subsection{Uniform geometry and fold coordinates}
We first verify the geometric facts needed to apply the local estimates
uniformly. At a simple zero, choose a product neighborhood where $g_s$
has fixed sign and is bounded away from zero. At a fold, choose one
where $g_{ss}$ has fixed sign and is bounded away from zero. A spatial
slice of the first neighborhood has at most one zero, and a slice of
the second has at most two, by strict monotonicity or strict convexity
(concavity). A finite cover of the compact zero set therefore bounds
the number of spatial zeros uniformly in $c$. On the closed part of
this set outside open fold neighborhoods, $|g_s|$ has a positive minimum.
Continuity also gives
\begin{equation}\label{eq:generic-anchor}
 \sup_{s,c} k_c(s)<\infty,\qquad
 \inf_{c\in I}\int_\Gamma k_c(s)\dd s>0.
\end{equation}
Indeed, a zero minimum of the latter integral would give an identically
zero realization, contradicting \cref{ass:generic-folds}.

For clarity, a fold coordinate can be constructed directly. Near
$(s_j,c_j)$ the implicit function theorem gives a critical point $z_j(c)$
with $g_s(z_j(c),c)=0$ and $z_j(c_j)=s_j$. Taylor's formula yields
\[
 g(s,c)=h_j(c)+(s-z_j(c))^2b_j(s,c),\qquad
 b_j(s,c)=\int_0^1(1-\theta)
 g_{ss}(z_j(c)+\theta(s-z_j(c)),c)\dd\theta,
\]
where $h_j(c)=g(z_j(c),c)$. After shrinking the neighborhood,
$b_j$ has fixed sign $\sigma_j$ and its absolute value is bounded
above and below by positive constants. The $C^2$ coordinate
$x=(s-z_j(c))\sqrt{|b_j(s,c)|}$ has a spatial Jacobian bounded
above and below, and
\begin{equation}\label{eq:generic-normal-form}
 g(s,c)=\sigma_j(x^2-t),\qquad t=-\sigma_j h_j(c),\qquad
 t(c_j)=0,\quad t'(c_j)=-\sigma_j g_c(s_j,c_j)\ne0.
\end{equation}
Thus $|t|\asymp|c-c_j|$, $|\dd c/\dd t|\asymp1$,
and $z_j(c)-s_j=O(|t|)$. No higher regularity of the normal form is used.

On the two-root side $t>0$, put $r=\sqrt t$. Each root has distance
comparable to $r$ from the fold site, their separation is comparable
to $r$, and their slopes have magnitude comparable to $r$. If $u$ is
either root, then on $|s-u|\le\eta r$, with one fixed small $\eta>0$,
\begin{equation}\label{eq:generic-root-potential}
 k_c(s)\asymp r^2(s-u)^2.
\end{equation}
Alternatively, $g_c\ne0$ writes the zero curve as $c=C_j(u)$, with
\[
 C_j'(s_j)=0,\qquad
 C_j''(s_j)=-\frac{g_{ss}(s_j,c_j)}{g_c(s_j,c_j)}\ne0.
\]
Hence $|C_j'(u)|\asymp|u-s_j|$. A root-position interval of length
comparable to $r$ at distance $r$ from the site has parameter
probability comparable to $r^2$.

Choose disjoint spatial neighborhoods of the fold sites and disjoint
parameter neighborhoods of the fold values. They may be shrunk so that
the unfolding roots remain inside the retained spatial neighborhood
whenever its parameter neighborhood is active. Outside these fold
products, the compact simple-zero set has uniformly nonzero slopes.
It therefore admits root intervals of a fixed small radius with
quadratic potential bounds and a uniform bound on their number.
These intervals can be taken disjoint after shrinking: otherwise a
sequence of distinct ordinary roots with separation tending to zero
would limit to an excluded multiple zero. On the remaining pieces,
outside the active fold neighborhood, the potential has a uniform
positive lower bound. These facts provide finite partitions for
Neumann bracketing. They impose no condition on $g_c$ at ordinary roots.

\subsection{Lower bounds from one sampled fold}
Fix one fold and one side of its zero curve. Set $s_j=0$ in a local
arclength coordinate. For $u\in[r,3r/2]$ let $I_r$ be a fixed
enlargement contained in $(r/2,2r)$, and let $m=\int_{I_r}D$.
The bounds on $C_j'$ and the density give
\[
 \varpi(C_j(u))|C_j'(u)|\asymp r
 \quad\hbox{for almost every }u\in[r,3r/2].
\]
Take the nonnegative compact bump used in \eqref{eq:blind-shell}.
If $0<m\le c_*r^7$, choose
$\ell=\delta_*(m/r^3)^{1/4}$, with fixed sufficiently small
constants $c_*$ and $\delta_*$.
All translated supports then lie in $I_r$ and satisfy
\eqref{eq:generic-root-potential}. Writing
\[
 B(u)=\ell^{-2}\int D(s)|\psi'((s-u)/\ell)|^2\dd s,
 \qquad \frac1r\int_r^{3r/2}B(u)\dd u\le\frac{Cm}{r\ell},
\]
the quotient formula and Jensen's inequality for the convex function
$x\mapsto x^{-q}$ give
\begin{align}
 \int_{C_j([r,3r/2])} J_c^g(D)^q\varpi(c)\dd c
 &\ge c_q r^2\ell^{2q}
       \left[\frac{m}{r\ell}+r^2\ell^3\right]^{-q}\notag\\
 &\ge c_q r^{2-5q/4}m^{-q/4}.
 \label{eq:generic-shell}
\end{align}
If $m=0$, the same bumps with widths tending to zero give an infinite
lower bound. This argument uses only local mass and remains valid for
arbitrarily concentrated or discontinuous mobility.

For $q<8/5$ one fixed shell and $m\le M$ give $c_qM^{-q/4}$.
At $q=8/5$, choose a geometric sequence of shells between a fixed
small radius and $CM^{1/7}$, with scale ratio large enough that both
their spatial enlargements and their parameter images are disjoint.
The constant $C$ ensures $m_n\le M\le c_*r_n^7$ for every retained
shell. There are $N\asymp\log(1/M)$ of them, and
\begin{equation}\label{eq:generic-critical-lower}
 \int_I J_c^g(D)^{8/5}\varpi(c)\dd c
 \ge c\sum_{n=1}^Nm_n^{-2/5}
 \ge cN^{7/5}M^{-2/5}.
\end{equation}
The last step uses $\sum m_n\le M$ and convexity of the negative power.

For $q>8/5$, put $R=M^{1/7}$ and use a bump of width $R$ centered
at $s_j$. For $|c-c_j|\le cR^2$, Taylor's formula gives
$k_c\le CR^4$ on its support. The derivative energy is at most
$CM/R^2$, the reaction energy at most $CR^5$, and the squared source
is comparable to $R^2$. Consequently $J_c^g(D)\ge cR^{-3}$
throughout a parameter window of probability at least $cR^2$.
This proves the lower order $R^{2-3q}=M^{(2-3q)/7}$.
Additional roots cannot weaken any of these local lower tests.

\subsection{A single design protecting all possible roots}
Let $\delta_F(s)$ be arclength distance to the finite set of fold
sites, and set
\begin{equation}\label{eq:generic-design}
 \alpha=\max\left\{0,\frac{6q-4}{q+4}\right\},\qquad
 D_M(s)=a_R(\delta_F(s)+R)^{-\alpha},\qquad
 a_R=\frac{M}{\int_\Gamma(\delta_F+R)^{-\alpha}\dd s}.
\end{equation}
Choose $R$ by \eqref{eq:graded-cutoffs}, replacing $\alpha_q$ by
$\alpha$ in its subcritical case. The normalization near the finitely
many sites gives
\begin{equation}\label{eq:generic-normalization}
 a_R\asymp
 \begin{cases}
 M,&\alpha<1,\\
 M/\log(1/M),&\alpha=1,\\
 MR^{\alpha-1},&\alpha>1,
 \end{cases}
 \qquad a_R\asymp R^{6+\alpha}.
\end{equation}
For each $M>0$ this field is bounded and bounded below by $ca_R>0$.
Using a nonnegative grading exponent matters: an ordinary root at
one parameter value can occupy a fold site associated with another
value. A negative exponent would reduce its local mobility. The
choice $\alpha=0$ for $q\le2/3$ protects such roots and suffices
for the stated orders.

Here are the uniform response estimates for this design. In a parameter
neighborhood of a fold, use the signed $t$ in
\eqref{eq:generic-normal-form}. There are fixed thresholds such that
\begin{align}
 J_c^g(D_M)&\le C\bigl[1+a_R^{-1/4}
              +a_R^{-1/4}r^{-3/2+\alpha/4}\bigr],
       &&t>0,\quad r=\sqrt t\ge CR,\label{eq:generic-root-upper}\\
 J_c^g(D_M)&\le C\bigl[1+a_R^{-1/4}+R^{-3}\bigr],
       &&|t|\le CR^2,\label{eq:generic-core-upper}\\
 J_c^g(D_M)&\le C\bigl[1+a_R^{-1/4}+|t|^{-3/2}\bigr],
       &&t<0,\quad |t|\ge CR^2.\label{eq:generic-out-upper}
\end{align}
The constant in the core threshold can be increased to cover the
transition between these regions. Away from all fold parameter
neighborhoods the bound is $C(1+a_R^{-1/4})$.

We justify these estimates, including their endpoint uniformity.
On an interval centered at a root, suppose the mobility is at least
$d$ and the potential is at least $a(s-u)^2$. If its half-width is
at least a fixed positive multiple of $\ell=(d/a)^{1/4}$,
rescaling and Neumann bracketing give
\begin{equation}\label{eq:generic-harmonic-upper}
 J_{\mathrm{interval}}^N\le C d^{-1/4}a^{-3/4}.
\end{equation}
The rescaled Neumann harmonic response is uniformly bounded as the
interval grows: retain a fixed central interval, whose derivative plus
quadratic reaction energy controls constants and hence has positive
coercivity, and integrate $x^{-2}$ on the two tails. Rescaled lengths
in a bounded range bounded away from zero have the same bound by
coercivity and compactness. This is the order version of
\cref{lem:interval-response,prop:harmonic}.

For a fold-side root with $r\ge CR$, take an interval of radius
$\eta r$. There $D_M\asymp a_Rr^{-\alpha}$ and
$k_c\asymp r^2(s-u)^2$, so its oscillator length divided by $r$ is
\[
 \asymp\left(\frac{a_R}{r^{6+\alpha}}\right)^{1/4}
 \asymp\left(\frac Rr\right)^{(6+\alpha)/4}.
\]
Equation \eqref{eq:generic-harmonic-upper} bounds the two root
contributions by $Ca_R^{-1/4}r^{-3/2+\alpha/4}$.
Outside these root intervals but within the fold neighborhood, the
normal form bounds the reciprocal-potential integral by $Cr^{-3}$:
scale $x$ by $r$ and exclude fixed neighborhoods of $x/r=\pm1$.
Its ratio to the root bound is at most
$C(R/r)^{(6+\alpha)/4}$, so it is absorbed.

In the central window $|t|\le CR^2$, bracket a spatial interval
$|s-s_j|\le KR$ with fixed sufficiently large $K$. On this interval
$D_M\asymp R^6$. With $s-s_j=Ry$ and $\tau=(c-c_j)/R^2$, the reaction
divided by $R^4$ differs uniformly by $o(1)$ from
$(\gamma_jy^2+\beta_j\tau)^2$, where
$\gamma_j=\tfrac12g_{ss}(s_j,c_j)$, $\beta_j=g_c(s_j,c_j)$, and
$\tau$ lies in a fixed compact interval. Both coefficients are nonzero.
The derivative term after rescaling has a uniform
positive coefficient. These interval forms have uniformly positive
Neumann coercivity: a sequence of unit $L^2$ functions with energy
tending to zero would converge to a nonzero constant, while a limiting
quadratic-square potential would force that constant to vanish.
Thus their source responses are $O(R^{-3})$. Choosing $K$ large
enough places all local roots inside the interval, and the remaining
local reciprocal-potential integral is also $O(R^{-3})$.

On the rootless side, bounded coordinate Jacobians in
\eqref{eq:generic-normal-form} give
\[
 \int_{\mathrm{fold\ neighborhood}}k_c^{-1}\dd s
 \le C\int_{\R}(|t|+x^2)^{-2}\dd x
 \le C|t|^{-3/2}.
\]
Every other root belongs to the uniform ordinary-root cover.
The global lower bound $D_M\ge ca_R$, its nonzero slope, and
\eqref{eq:generic-harmonic-upper} bound each such contribution by
$Ca_R^{-1/4}$, even if that root is stationary in $c$.
The remaining positive-potential pieces contribute a constant.
Allowing independent endpoint traces enlarges the variational class,
so summing these interval responses bounds the full response from
above. The number of intervals is uniformly bounded. This proves
\cref{eq:generic-root-upper,eq:generic-core-upper,eq:generic-out-upper}.

To finish the moment estimate, a finite sum of nonnegative responses
raised to a fixed $q>0$ is bounded by a fixed multiple of the sum of
their $q$th powers. The bounded density and $|\dd c/\dd t|\asymp1$
reduce the root-side contribution to
\begin{equation}\label{eq:generic-moment-integral}
 Ca_R^{-q/4}\int_R^{r_*}r^{e-1}\dd r,\qquad
 e=2-\frac{3q}{2}+\frac{\alpha q}{4}
 =\begin{cases}
 2-3q/2,&q\le2/3,\\
 (8-5q)/(q+4),&q>2/3.
 \end{cases}
\end{equation}
The ordinary-root contribution is $O(a_R^{-q/4})$; the central
window contributes $O(R^{2-3q})$. The rootless integral is bounded
when $q<2/3$, logarithmic when $q=2/3$, and $O(R^{2-3q})$ when
$q>2/3$.

For $q<8/5$, $e>0$ and $a_R\asymp M$, so the root integral is
$O(M^{-q/4})$. The central power satisfies
$R^{2-3q}/a_R^{-q/4}\asymp R^e\to0$; any bounded or logarithmic
rootless term is smaller as well. At $q=8/5$, $\alpha=1$ and $e=0$,
and the root integral is
\[
 O\bigl(a_R^{-2/5}\log(1/R)\bigr)
 =O\bigl(M^{-2/5}[\log(1/M)]^{7/5}\bigr).
\]
The central and rootless terms have one fewer logarithmic factor.
For $q>8/5$, $e<0$ and
$a_R^{-q/4}R^e\asymp R^{2-3q}=M^{(2-3q)/7}$;
the regular contribution is smaller. These are the upper bounds in
\cref{thm:generic-orders}, completing its proof.

\subsection{Consequences for the physical transport problem}
Let $\mathcal G_q^g(M)$ be the corresponding optimal full-bulk moment
over physical predetermined fields, with the fixed geometry, positive
bulk diffusivity, constant affinity, and nonzero $V$ assumed in
\cref{sec:model}. The uniform bounds \eqref{eq:generic-anchor} verify
\eqref{eq:rate-assumptions}. All three scales in
\eqref{eq:generic-three} dominate $[1+\log(1/M)]^q$. Therefore
\cref{thm:optimal-transfer}, with no observation, gives the stronger
comparison
\begin{equation}\label{eq:generic-bulk}
 \mathcal G_q^g(M)\sim\chi^q\mathcal P_q^g(M).
\end{equation}
In particular the same three moment orders hold for finite transverse
bulk diffusion. The exact ratio follows by adding a vanishing uniform
background to a scalar near-minimizer at the same total budget, as
in \eqref{eq:mixture}; it does not require that the explicit graded
trial be a sharp optimizer. Directly on that trial,
$\min D_M\ge ca_R$ and
$\log(1/a_R)=O_q(\log(1/M))$, so \cref{thm:log-remainder} bounds
the bulk correction by a lower-order logarithm. No bounded correction
or correction limit is asserted for a general family.

Only one positively sampled fold is used in the lower bounds; the
upper design protects every possible fold site. Thus asymmetry, unequal
root slopes, and additional regular roots change comparison constants
without changing the threshold. In contrast, vanishing or singular
sampling density at a fold, higher spatial degeneracy, failure of
$g_c\ne0$, or a family approaching an identically zero profile falls
outside the proof. A law supported only on the rootless side need not
have the moving-root contribution. Mobility caps, prescribed positive
backgrounds, fixed spatial resolution, and uncertainty in the fold
sites are also different design problems. The result uses the known
fold geometry and an unrestricted integral budget; it does not claim
generic sharp coefficients or a unique limiting design.
